% ══════════════════════════════════════════════
% CHƯƠNG 2 — PHÂN TÍCH VÀ THIẾT KẾ CƠ SỞ DỮ LIỆU
% Người viết: TV1 (Lead + DB)
% ══════════════════════════════════════════════

\chapter{Phân tích và Thiết kế Cơ sở dữ liệu}

\section{Sơ đồ thực thể mối quan hệ (ERD)}\label{Sec2.1}

\begin{de}\label{def2.1}
  % [TV1] Mô tả tổng quan cơ sở dữ liệu: có bao nhiêu bảng,
  % được tổ chức thành mấy nhóm chức năng, tên từng nhóm là gì?

  % [Điền nội dung vào đây]
\end{de}

% [TV1] Xuất ERD từ SSMS hoặc dbdiagram.io, lưu thành file erd.png
%       rồi thay dòng \framebox bên dưới bằng: \includegraphics[width=14cm]{erd.png}
\begin{figure}[H]
  \centering
  \framebox[14cm][c]{\rule{0pt}{8cm} [Chèn hình ERD vào đây]}
  \caption{Sơ đồ thực thể mối quan hệ (ERD) của hệ thống}\label{fig:erd}
\end{figure}

\section{Mô tả các bảng chính}\label{Sec2.2}

% [TV1] Điền tên bảng, số cột và mô tả chức năng của từng bảng.
% Thêm/bớt dòng tùy theo số bảng thực tế của nhóm.

\begin{table}[H]
  \begin{center}
    \begin{tabular}{|l|c|l|}
      \hline
      \textbf{Tên bảng} & \textbf{Số cột} & \textbf{Mô tả} \\
      \hline
      % [TênBảng1] & [n] & [Mô tả chức năng bảng] \\
      \hline
      % [TênBảng2] & [n] & [Mô tả chức năng bảng] \\
      \hline
      % [TênBảng3] & [n] & [Mô tả chức năng bảng] \\
      \hline
      % [TênBảng4] & [n] & [Mô tả chức năng bảng] \\
      \hline
      % [TênBảng5] & [n] & [Mô tả chức năng bảng] \\
      \hline
      % ... thêm dòng nếu cần
    \end{tabular}
  \end{center}
  \caption{Mô tả các bảng trong cơ sở dữ liệu}\label{table:tables}
\end{table}

\section{Các ràng buộc dữ liệu}\label{Sec2.3}

\begin{lem}\label{lem2.1}
  % [TV1] Phát biểu tổng quát về việc áp dụng ràng buộc dữ liệu
  % (CHECK constraints, NOT NULL, UNIQUE, FK...) tại tầng CSDL.

  % [Điền nội dung vào đây]
\end{lem}

% [TV1] Liệt kê các ràng buộc quan trọng nhóm đã thiết kế.
% Dùng cú pháp \texttt{} cho tên cột/bảng và giá trị kỹ thuật.

\begin{itemize}
  \item \texttt{} % [Ràng buộc 1 — ví dụ: TênBảng.TênCột > 0]
  \item \texttt{} % [Ràng buộc 2]
  \item \texttt{} % [Ràng buộc 3]
  \item \texttt{} % [Ràng buộc 4]
  \item \texttt{} % [Ràng buộc 5 — thêm nếu cần]
\end{itemize}

\section{Views}\label{Sec2.4}

\begin{thm}\label{thm2.1}
  % [TV1] Phát biểu: nhóm tạo bao nhiêu View, mục đích tổng quát là gì?

  % [Điền nội dung vào đây]
\end{thm}

% [TV1] Với mỗi View: viết 1 dòng mô tả rồi chèn code SQL vào lstlisting.
% Sao chép cấu trúc bên dưới cho View 2, View 3, ...

\noindent\textbf{View 1 — [Tên View]:} % [Mô tả ngắn View 1]

\begin{lstlisting}[style=sql, caption={View [TênView1]}]
-- [TV1] Dán code SQL CREATE VIEW vào đây
\end{lstlisting}

\noindent\textbf{View 2 — [Tên View]:} % [Mô tả ngắn View 2]

\begin{lstlisting}[style=sql, caption={View [TênView2]}]
-- [TV1] Dán code SQL CREATE VIEW vào đây
\end{lstlisting}

\noindent\textbf{View 3 — [Tên View]:} % [Mô tả ngắn View 3]

\begin{lstlisting}[style=sql, caption={View [TênView3]}]
-- [TV1] Dán code SQL CREATE VIEW vào đây
\end{lstlisting}

\section{Functions}\label{Sec2.5}

% [TV1] Mô tả tên và chức năng của từng Function T-SQL.

\noindent Các Function T-SQL được xây dựng:

\begin{enumerate}
  \item \textbf{[TênFunction1]}: % [Mô tả chức năng]
  \item \textbf{[TênFunction2]}: % [Mô tả chức năng]
  \item \textbf{[TênFunction3]}: % [Mô tả chức năng — thêm nếu có]
\end{enumerate}

% [TV1] Chèn code SQL của ít nhất 1 Function tiêu biểu.

\begin{lstlisting}[style=sql, caption={Function [TênFunction]}]
-- [TV1] Dán code SQL CREATE FUNCTION vào đây
\end{lstlisting}

\section{Stored Procedures}\label{Sec2.6}

\begin{thm}\label{thm2.2}
  % [TV1] Phát biểu về Stored Procedure quan trọng nhất (ví dụ: dùng
  % Transaction + lock để đảm bảo ACID khi nhiều người thao tác đồng thời).

  % [Điền nội dung vào đây]
\end{thm}

\noindent\begin{proof}
  % [TV1] Giải thích tại sao thiết kế của SP đảm bảo tính đúng đắn
  % (ví dụ: cơ chế khóa dòng, TRY-CATCH, ROLLBACK khi có lỗi...).

  % [Điền nội dung vào đây]
  \QEDAa
\end{proof}

% [TV1] Chèn code SQL đầy đủ của từng Stored Procedure.
% Sao chép cấu trúc lstlisting bên dưới cho mỗi SP.

\begin{lstlisting}[style=sql, caption={Stored Procedure [TênSP1]}]
-- [TV1] Dán code SQL CREATE PROCEDURE vào đây
\end{lstlisting}

\begin{lstlisting}[style=sql, caption={Stored Procedure [TênSP2]}]
-- [TV1] Dán code SQL CREATE PROCEDURE vào đây
\end{lstlisting}

% Thêm lstlisting cho SP3, SP4 nếu có

\section{Triggers}\label{Sec2.7}

\begin{ex}\label{ex2.1}
  % [TV1] Mô tả Trigger đầu tiên: tên, sự kiện kích hoạt (AFTER INSERT/UPDATE/DELETE),
  % và tác dụng của nó (tự động cập nhật bảng nào, trường nào?).

  % [Điền nội dung vào đây]
\end{ex}

\begin{lstlisting}[style=sql, caption={Trigger [TênTrigger1]}]
-- [TV1] Dán code SQL CREATE TRIGGER vào đây
\end{lstlisting}

\begin{cor}
  % [TV1] Nêu hệ quả/lợi ích của Trigger vừa trình bày
  % (ví dụ: tách trách nhiệm, tránh xử lý trùng lặp...).

  % [Điền nội dung vào đây]
\end{cor}

% [TV1] Thêm lstlisting cho Trigger thứ 2 nếu có

\section{Kết quả kiểm thử T-SQL}\label{Sec2.8}

% [TV1] Chạy các test script trong SSMS, chụp màn hình kết quả
%       và lưu thành file (ví dụ: ssms_test.png).
%       Sau đó thay \framebox bên dưới bằng: \includegraphics[width=14cm]{ssms_test.png}

\begin{figure}[H]
  \centering
  \framebox[14cm][c]{\rule{0pt}{5cm} [Chèn screenshot kết quả test SSMS]}
  \caption{Kết quả kiểm thử Stored Procedure trên SSMS}\label{fig:ssms}
\end{figure}
